\section{Theoretical background}
\label{sec:background}

\subsection{Misinformation as exposure, choice, and diffusion}

Research on online misinformation distinguishes production, exposure, belief, and sharing. These events are related but not identical. The existence of false content does not imply that every user is exposed to it, and exposure does not imply belief or reposting. Empirical studies show that false-news consumption and sharing can be concentrated among relatively small groups \citep{grinberg2019,guess2019}. At the same time, false stories can travel farther and faster than truthful stories under particular platform and content conditions \citep{vosoughi2018}. This combination motivates metrics that retain both the source of a decision and the truth state of the item selected.

Source credibility offers a useful but incomplete heuristic. Crowdsourced source-quality judgments can help identify unreliable outlets \citep{pennycook2019}, and accuracy prompts can redirect attention toward truthfulness \citep{pennycook2020,epstein2021}. Yet source labels and reputations operate alongside content-level features. Emotion, novelty, simplicity, proximity, sensationalism, multimedia, and other engagement cues may shape selection before a user knows whether the item is true. Moreover, a source can alter its style without altering its underlying editorial behaviour. A model that collapses perceived credibility and actual truth generation into one variable risks obscuring this difference.

The wider information environment also matters. Algorithmic and individual choices jointly structure exposure \citep{bakshy2015}. Homogeneous communities can reinforce selective narratives \citep{delvicario2016}, while social influence can bias evaluations even when the initial social signal is arbitrary \citep{muchnik2013}. Network experiments further show that diffusion depends on reinforcement structure and the complexity of the behaviour being transmitted \citep{centola2010}. These findings do not determine a unique simulation mechanism, but they establish that content attributes, social transmission, and network opportunity should be represented separately enough to support diagnosis.

\subsection{Endorsements as bounded evidence}

Endorsement theory represents reasoning under uncertainty through labelled items of evidence whose interpretation depends on their source, recency, and relevance \citep{cohen1985}. Instead of assigning a single immutable utility to an alternative, an agent accumulates endorsements and evaluates them in context. This is attractive for social simulation because users can hold heterogeneous weights, observe incomplete information, and update source evaluations over time.

The prior e-commerce model by \citet{teran2022} uses endorsements to represent marketplace choice and WOM. Buyers evaluate marketplace attributes, retain evidence, and communicate choices. \citet{teran2024} extends that work through scenario analysis and repeated Monte Carlo experiments, emphasising how recognition, attributes, and WOM shape market share. FAKENEWS-ABM preserves the reasoning architecture while changing the domain interpretation. News sources replace marketplaces; users replace buyers; a selected source denotes a repost; and source-level truth states make it possible to calculate realised false reposts.

This transfer creates an important conceptual caution. Market share is a suitable outcome when the primary phenomenon is marketplace dominance. In misinformation research, source share is only one layer. A source can dominate while publishing mostly true content, or remain relatively small while producing a high proportion of false items. Accordingly, the present study treats target-source selection as an intermediate outcome and actual fake-repost share as the primary dissemination outcome.

\subsection{WOM, truth feedback, and memory}

WOM is interpersonal transmission about an alternative. In the model it performs two logically distinct functions. First, a recommendation can make an unknown source available for later evaluation. Second, after the recommended publication's truth state is known, the receiver can add a signed endorsement to that source. The first function is discovery; the second is truth-sensitive social learning. Keeping them separate allows a baseline test: if discovery-only WOM and no WOM are identical, any difference created by truth-sensitive WOM can be attributed within the model to the signed outcome rule rather than to source discovery alone.

Signed feedback is configurable. A false publication can be penalised ($-1$), ignored ($0$), or rewarded ($+1$); a true publication can be treated in the same way, although the current experiments focus on ignoring or rewarding true outcomes. Ignoring means that no outcome endorsement is created. It does not store a zero-valued endorsement that might later be recoded as negative. This distinction is necessary because a binary conversion of zero to ``not positive'' would silently turn neutrality into punishment.

Memory determines how long accumulated endorsements remain active. Short memory creates rapid turnover: copied attributes can influence only recent observations, and an intervention needs repeated reinforcement to persist. Long or infinite memory allows early evidence to remain in the decision set. The resulting effects need not be monotonic. Long memory can preserve treatment-generated endorsements, but it can also preserve pre-intervention evidence or stabilise a control. Therefore memory must be analysed jointly with activation time and WOM policy rather than interpreted as a simple scalar measure of susceptibility.

\subsection{Scenario analysis and simulation evidence}

Agent-based models are especially useful when aggregate outcomes arise from heterogeneous local decisions and feedback \citep{macal2010,railsback2019}. Their credibility depends on transparent specification, verification, validation, and disciplined experimentation. The ODD protocol encourages explicit descriptions of purpose, entities, state variables, process scheduling, and submodels \citep{grimm2020}. Monte Carlo replication separates persistent tendencies from individual stochastic realisations, and common random numbers can improve contrasts when treatment and control respond similarly to shared randomness \citep{law2015}.

Scenario analysis does not predict a single future. It asks conditional questions: if a source copied a specified set of attributes at a specified time, and if users retained evidence for a specified horizon, what would the implemented mechanisms produce? The strength of the answer depends on internal validity and sensitivity analysis. External claims require empirical calibration. In this project, a calibration interface exists, but observed repost, exposure, credibility, and uncertainty targets are not part of the repository. We therefore report mechanism-based simulation contrasts and state empirical calibration as future work.
